% =============================================================================
% Appendix: dataset revision (removal of questions with unrecoverable defects).
%
% Usage: \input this file after \appendix in the paper.
% Required packages: booktabs. Assumes UTF-8 input and T1 font encoding (both
% provided by the ACL template). No figures, no bibliography.
%
% Every number comes from config/removed_questions.json, the revised outputs
% output/04 - revised and output/06 - revised, the abstention subset
% output/knowledge-annotation-work/subsets/full-run-2-uncertain-revised, the
% exports and manifests of the annotation runs full-run-2,
% uncertain-taxonomy-v1.3, uncertain-taxonomy-v1.4 and full-run-taxonomy-v1.4
% (with its post-annotation/ directory) and the audit
% taxonomy-v1.2-audit/reviewed_cases.json, and is printed by
% ablations/revision/make_summary.py, which also checks the artifacts against
% each other. The knowledge-area appendix is mentioned in prose only and is
% never referenced with \ref, so this file compiles on its own.
% =============================================================================

\section{Dataset Revision: Removal of Questions with Unrecoverable Defects}
\label{app:rev}

Nineteen questions were removed from the filtered corpus of 41,653 questions
(0.046\%) because the task they pose cannot be recovered from the dataset: the
statement was replaced or truncated during extraction or pre-processing, the
supporting text, chart or table that the statement refers to was lost, or the
item is one of a block that can only be solved jointly. As evaluation items
they could only be answered by chance. None of them was removed for being
hard, for receiving an unexpected label or because the annotation passes
disagreed. The defects were found as a by-product of the knowledge-area
annotation, whose annotator may abstain, and completed by a string-pattern
sweep. This appendix describes the defect families and the removal criterion
(\S\ref{app:rev-criteria}), how the questions were found
(\S\ref{app:rev-identification}), the edge cases that were read and kept
(\S\ref{app:rev-validation}), the procedure that applied the removal with an
audit trail (\S\ref{app:rev-procedure}), its effect on the published datasets
(\S\ref{app:rev-effect}) and its limitations (\S\ref{app:rev-limitations}).

\paragraph{Summary.} (1)~Five defect families, 19 questions from six exams:
six self-referential OBI items, seven BLUEX items whose statement was
replaced by a formatting instruction, two truncated POSCOMP items, two
identical IME items without their reading passage and two items (COMVEST,
ENEM) whose chart or table is missing. (2)~Ten were abstentions of the
annotation run with taxonomy~1.2, seven were found by searching the corpus
for the same defect patterns, and two were abstentions of the final run.
(3)~Three groups of look-alike questions were read and kept: 78 of the 80
questions whose statement starts with an ``INSTRUÇÕES'' header, 19 questions
whose last choice carries the heading of the next reading passage, and two
questions whose supporting text is missing but whose task is identifiable.
(4)~Removal is applied by a content hash of exam, edition, item number,
statement and choices, never by row index, into a new directory with a
manifest, a removal log and a reload check; the original stages are never
modified. (5)~The filtered corpus goes from 41,653 to 41,634 questions and
the deduplicated corpus from 37,887 to 37,873; five of the 19 questions had
already been removed by the deduplication. (6)~The three affected Hub
datasets were republished at revisions \texttt{e0f638e0}\ldots\ (revised
corpus), \texttt{54ba11ff}\ldots\ (deduplicated corpus) and
\texttt{fb9d509a}\ldots\ (annotated corpus).

% -----------------------------------------------------------------------------
\subsection{Defect families and removal criterion}
\label{app:rev-criteria}

Table~\ref{tab:rev-removed} lists the 19 questions by family.

\paragraph{Self-referential block.} The second phase of the 2012 Brazilian
Informatics Olympiad opens with three items whose statement is only ``A
resposta da questão $N$ é:'' (``The answer to question $N$ is:'') and whose
choices are the letters A to E; item~1 refers to item~2, item~2 to item~3 and
item~3 to item~1. The three form one logic puzzle that is solved jointly. As
independent rows, which is what the dataset format preserves, each is
undecidable. The block appears in the level-1 and the level-2 booklets, hence
six rows.

\paragraph{Statement replaced by a formatting instruction.} In seven BLUEX
items (UNICAMP 2018, 2020, 2023 and twice 2024; USP 2020 and 2022) the
statement is the Portuguese instruction ``Seu único objetivo é detectar as
expressões matemáticas, convertê-las para o padrão LaTeX\ldots'' (``Your only
goal is to detect the mathematical expressions, convert them to LaTeX\ldots'')
followed by a worked example about a function $h(x)$. The instruction belongs
to the pre-processing of the source and overwrote the original statement. The
choices are those of the original literature, history or grammar question and
are orphaned; the original statement is not recoverable from the dataset.

\paragraph{Truncated statement.} Two POSCOMP 2024 items (19 and 20) kept only
the header ``INSTRUÇÕES'' and the final question (``the percentage of time in
which the computation takes less than 65 seconds is'', ``the approximate mean
time the software takes to compute is''), without the data or the algorithm
that define the computation time.

\paragraph{Missing supporting text.} Item~36 of the IME entrance exams of
2010 and 2011 is the English-comprehension question ``What task below could
Lammert B. Otten be legally in charge of?'', with the choices ``suing a
criminal'', ``prescribing drugs'', ``piloting a plane'', ``saying a mass'' and
``filling in a cavity''. The reading passage that introduces the person is
absent. The two rows are identical in statement and choices and carry
different answer keys (A in 2010, D in 2011).

\paragraph{Missing supporting data.} COMVEST 2017 item~63 refers to ``o
gráfico abaixo'' (``the chart below'') and is cut before the relation or the
values it asks about; ENEM 2018 item~136 is the single sentence ``Com base
nessas informações, o banco que transferiu a maior quantia via TED é o banco''
(``Based on this information, the bank that transferred the largest amount is
bank'') with the choices 1 to 5 and no table.

\paragraph{Criterion.} A question is removed on one of two grounds only: its
statement or supporting material cannot be recovered from the dataset, so
that the task is not answerable; or it depends on a block of items that the
one-row-per-question format does not preserve. Questions whose supporting
material is missing but whose task is still identifiable from the statement
and the choices are kept, and treated as an annotation matter
(\S\ref{app:rev-validation}). Difficulty, disagreement between the annotation
passes and unexpected labels are not grounds for removal.

\begin{table*}[t]
\centering
\footnotesize
\setlength{\tabcolsep}{4pt}
\begin{tabular}{@{}lp{6.6cm}rrp{4.2cm}@{}}
\toprule
Family & Exam, edition and item & Stage 04 & Stage 06 & Detection \\
\midrule
Self-referential block & OBI 2012, phase 2, levels 1 and 2, items 1, 2 and 3 & 6 & 3 & abstention, run 1.2 (6) \\
Statement replaced by instruction & BLUEX: UNICAMP 2018 item 2; UNICAMP 2020 item 68; UNICAMP 2023 item 9; UNICAMP 2024 items 7 and 12; USP 2020 item 37; USP 2022 item 65 & 7 & 6 & abstention, run 1.2 (1); pattern sweep (6) \\
Truncated statement & POSCOMP 2024 items 19 and 20 & 2 & 2 & abstention, run 1.2 (1); pattern sweep (1) \\
Missing supporting text & IME 2010 and 2011, item 36 & 2 & 1 & abstention, run 1.2 (2); confirmed by run 1.4 on the abstention subset \\
Missing supporting data & COMVEST 2017 item 63; ENEM 2018 item 136 & 2 & 2 & abstention, run 1.4 (2) \\
\midrule
Total & & 19 & 14 & \\
\bottomrule
\end{tabular}
\caption{The removed questions. \emph{Stage 04}: rows removed from the filtered corpus (41,653 questions); \emph{Stage 06}: rows removed from the semantically deduplicated corpus (37,887 questions), where five of the listed questions were already absent because the exact deduplication had removed them as copies of another listed question (the level-2 copies of the three OBI items, UNICAMP 2024 item 12 and the 2010 copy of the IME item). \emph{Detection}: how each question entered the list, with the number of questions in parentheses; run 1.2 and run 1.4 are the full knowledge-area annotation runs with taxonomies 1.2 and 1.4.}
\label{tab:rev-removed}
\end{table*}

% -----------------------------------------------------------------------------
\subsection{Identification}
\label{app:rev-identification}

The questions were not found by a dedicated quality audit of the corpus but
through the knowledge-area annotation, in which a language model assigns a
subject to every question and may abstain with \texttt{UNCERTAIN}. A
statement that does not pose a task tends to produce an abstention, so the
abstentions of each full run were read, and the ones whose cause was the
source rather than the taxonomy were set aside. Table~\ref{tab:rev-funnel}
follows the list through the steps described next.

\paragraph{Abstentions as a defect detector.} The engineering audit of the
first full run (taxonomy~1.1) read 140 of its labels and abstentions and
tagged 23 of them as source issues; 15 of the 19 questions removed here are
among those 23, but at that point they were only marked for correction at the
source. The second full run, with taxonomy~1.2, ended with 23 abstentions on
the 41,653 questions. Running the experimental taxonomy~1.3 on those 23
questions labelled 14 and left 9 abstentions. The 9 were read: 8 had no
recoverable statement (the six OBI items, UNICAMP 2018 item~2 and POSCOMP 2024
item~20) and the ninth, AFA 2019 item~45, is an English-comprehension item
without its passage but with an identifiable task, which was kept.

\paragraph{Pattern sweep.} The 8 questions were then used as patterns over
the filtered and the deduplicated stages. The literal instruction string
matches exactly 7 rows of the filtered corpus: UNICAMP 2018 item~2 and 6
BLUEX items that run~1.2 had labelled from their orphaned choices, as
Portuguese Language, Literature, History or, in one case, Mathematics, with
confidences between 0.85 and 0.95; the string ``A resposta da questão''
matches exactly the 6 OBI rows; and the ``INSTRUÇÕES'' header matches 80
rows, among which only POSCOMP 2024 item~19 shares the truncation of item~20
(\S\ref{app:rev-validation}). The sweep added 7 questions that the annotator
had not flagged, and the list reached 15 questions on 7~October~2026. The
abstention subset of run~1.2 was revised with that list (23 to 15 questions)
and became the input of the experiment with taxonomy~1.4.

\paragraph{Identical copies with different labels.} Run~1.4 on the 15
questions labelled 14 and abstained on one, the 2011 copy of IME item~36,
while giving English Language to its identical 2010 copy; in run~1.2 both had
been abstentions and in run~1.3 both had received English Language. Two rows
with the same statement and choices oscillating between a label and an
abstention confirmed that the question depends on the missing passage, and
both were removed, bringing the list to 17. The filtered corpus revised with
those 17 questions (41,636 rows) was published as the base of the third full
run.

\paragraph{Abstentions of the final run.} The third full run, with
taxonomy~1.4, ended with 8 abstentions, none of which had been an abstention
of an earlier run. All 8 were read: 6 received a subject by hand, and 2 were
removed on 8~October~2026 because the chart or table their statement refers to
is missing (COMVEST 2017 item~63, ENEM 2018 item~136). The list reached its
final size of 19, and the revised filtered and deduplicated corpora were
regenerated from the original stages with the complete list.

\begin{table}[t]
\centering
\footnotesize
\setlength{\tabcolsep}{4pt}
\begin{tabular}{@{}lrrrr@{}}
\toprule
Step & Examined & Unc. & Added & List \\
\midrule
Full run 1.2 & 41,653 & 23 & -- & 0 \\
Run 1.3 on the abstentions & 23 & 9 & -- & 0 \\
Reading of the 9 abstentions & 9 & 1 & 8 & 8 \\
Pattern sweep, stages 04 and 06 & 41,653 & -- & 7 & 15 \\
Run 1.4 on the abstentions & 15 & 1 & -- & 15 \\
Reading of the IME copies & 2 & -- & 2 & 17 \\
Full run 1.4 & 41,636 & 8 & -- & 17 \\
Post-annotation & 8 & 0 & 2 & 19 \\
\bottomrule
\end{tabular}
\caption{How the removal list was built. \emph{Examined}: questions annotated or read at the step; \emph{Unc.}: questions left \texttt{UNCERTAIN} by the step (for the reading steps, abstentions kept in the dataset); \emph{Added}: questions added to the list; \emph{List}: size of the list after the step. The list was created on 7~October~2026 with 15 questions and completed on 8~October~2026. The 6 abstentions of the post-annotation that were not removed received a subject by hand.}
\label{tab:rev-funnel}
\end{table}

% -----------------------------------------------------------------------------
\subsection{Manual validation of edge cases}
\label{app:rev-validation}

Before the removal was applied, the groups of questions that share a surface
feature with a removed question were read one by one and kept, because they
do not share the defect.

\paragraph{Questions with an instruction header.} 80 rows start with
``INSTRUÇÕES'' or ``Instruções:'' (65 from POSCOMP 2024, 13 from three BACEN
editions, two of 2005 and 2006 and one undated, one from ENADE 2006 and one
from ENADE 2007). In all but the two POSCOMP items removed, the full
statement follows the header, including the shared supporting text of the
BACEN and ENADE question blocks. POSCOMP 2024 item~66,
whose statement is only ``Assinale a alternativa correta'' (``Mark the correct
alternative''), was kept because its choices are self-contained assertions
about the IP protocol.

\paragraph{Next passage appended to the last choice.} In 19 rows of COMVEST
2011 to 2015 (7, 5, 2 and 3 rows in 2011, 2012, 2014 and 2015) and of BLUEX
USP 2020 and 2024 (one each), the last choice carries, after its real text,
the heading of the next reading passage (``Texto para as questões\ldots''). The
correct choice and the others are intact before the noise, so the item
remains answerable and identifiable. These rows are recorded as a formatting
clean-up to be made, not as removals.

\paragraph{Missing support with an identifiable task.} AFA 2019 item~45 (an
English fragment with English choices) and COMVEST 2015 item~38 (a technique
used by Mendel, with its infographic missing) were kept: the competence being
assessed is recognizable from the statement and the choices, and the question
of which subject to assign is handled by the annotation taxonomy. The two IME
items were initially in this group and moved to the removal list only after
run~1.4 split their labels.

% -----------------------------------------------------------------------------
\subsection{Procedure}
\label{app:rev-procedure}

\paragraph{Identity.} Each listed question is identified by the SHA-256 of
the canonical JSON of its exam, edition, item number, statement and choices,
the same identifier the annotation uses. Removal is applied by identity, never
by row index, so one list serves every pipeline stage and every subset, whose
row orders differ. The versioned list records, for each question, the
identifier, exam, edition and item number, a category, the reason and the
evidence that brought it to the list, together with the creation and update
dates.

\paragraph{Tool.} The \texttt{revise} command of \texttt{mmlu\_pt.revision}
reads a dataset saved to disk, refuses to write into an existing directory and
never modifies its input. It writes the revised dataset, the same records as
JSON Lines, a removal log with the input row index, identifier, exam,
edition, item number, category and reason of every removed row, and a
manifest with the SHA-256 of the input files, the output files and the list,
the number of input, removed and output rows, and the listed identifiers that
were absent from the input. If the exam, edition or item number of a row
differ from the list entry with the same identifier, the run aborts. After
writing, the tool reloads the output and verifies that the kept rows are
identical to the input rows and in the same order before publishing the
directory. The \texttt{publish} command re-hashes the output before pushing
it to the Hub and records the repository and the resulting revision in the
manifest. The last two removals were also applied to the export of run~1.4 by
the post-annotation command of the annotation tool, under the same list and
identity, so that the annotated dataset and the revised corpus contain the
same identifiers in the same order.

% -----------------------------------------------------------------------------
\subsection{Effect on the datasets}
\label{app:rev-effect}

Table~\ref{tab:rev-datasets} gives the row counts. The filtered corpus
(stage~04 of the pipeline, published as \texttt{bench-temp/mmlu-pt-filtered})
loses all 19 questions and has 41,634 rows. The semantically deduplicated
corpus (stage~06) loses 14: the other five listed questions were already
absent because the exact deduplication, whose key is the normalized
statement, had removed each of them as a copy of another listed question that
the stage kept (the level-2 copies of the three OBI items, identical to the
level-1 copies; UNICAMP 2024 item~12, whose statement is identical to that of
UNICAMP 2018 item~2; and the 2010 copy of the IME item). By exam,
the removal takes 7 questions from BLUEX, 6 from OBI, 2 from POSCOMP, 2 from
IME and one each from COMVEST and ENEM. The intermediate revision with 17
questions (41,636 rows) was the input of the third full annotation run and is
superseded; it survives only as the source revision recorded in that run's
manifest. The identifiers of the final annotated dataset coincide, in the
same order, with those of the revised filtered corpus.

\begin{table*}[t]
\centering
\footnotesize
\setlength{\tabcolsep}{4pt}
\begin{tabular}{@{}lrrrrll@{}}
\toprule
Dataset & Input & Removed & Absent & Output & Hub repository & Revision \\
\midrule
Filtered corpus, 17-question list & 41,653 & 17 & 0 & 41,636 & \texttt{bench-temp-2/mmlu-pt-revised} & \texttt{0e8b69a1}\ldots\ (superseded) \\
Filtered corpus, 19-question list & 41,653 & 19 & 0 & 41,634 & \texttt{bench-temp-2/mmlu-pt-revised} & \texttt{e0f638e0}\ldots \\
Deduplicated corpus & 37,887 & 14 & 5 & 37,873 & \texttt{bench-temp-2/mmlu-pt-deduplicated} & \texttt{54ba11ff}\ldots \\
Abstentions of run 1.2, 15-question list & 23 & 8 & 7 & 15 & -- & -- \\
Annotated export of run 1.4 & 41,636 & 2 & 0 & 41,634 & \texttt{bench-temp-2/mmlu-pt-knowledge-annotated} & \texttt{fb9d509a}\ldots \\
\bottomrule
\end{tabular}
\caption{Datasets revised with the list, in the order in which they were produced. \emph{Absent}: listed questions not present in the input, so not removed from it; for the deduplicated corpus they had been removed by the deduplication, and for the abstention subset they were never abstentions. The abstention subset is a local working dataset and was not published. The three published revisions were pushed on 8~October~2026.}
\label{tab:rev-datasets}
\end{table*}

% -----------------------------------------------------------------------------
\subsection{Limitations}
\label{app:rev-limitations}

The detector is the abstention of one annotation model, complemented by
string patterns derived from the abstentions it produced; only defects that
make the model abstain or that share a pattern with one of them were found.
The corpus was not audited question by question by a person, and the sweep is
string-based and tied to the exams in which the patterns occurred, so
truncated statements that lack a distinctive marker, such as the two POSCOMP
items, are found only when the annotator abstains. The audit of the first
full run tagged 23 questions as source issues and 8 of them remain in the
dataset: six whose photograph, chart or text is missing but whose task was
judged identifiable (ENADE 2004 item~25 of the occupational-therapy booklet,
four BNDES items of 2008 and 2010, AFA 2019 item~45) and two that were not
re-examined under the removal criterion, BLUEX UNICAMP 2023 item~3, whose
statement field contains an explanatory text rather than a question, and ENEM
2020 item~1, whose poster is missing. The 19 rows with the next passage
appended to a choice, and other extraction artefacts of the same kind, are
still in the dataset. ``Unrecoverable'' was judged against the dataset, not
against the source PDFs, from which a statement could in principle be
re-extracted. Finally, the identical IME copies were caught only because the
annotation diverged between them; copies that received the same label would
have escaped the detector.
